\section{Reading a boundary value with residues}\label{sec:cartier}

Proposition~\ref{prop:descent-curve} gives a curve on which the output
coefficients $Y(X)$ have the required value $c$ at a place above $t=0$.
Individual parameter monomials can have poles there. For instance,
if $u=t^{-1}$ and $v=t^{-1}-t$, then $u^2$ and $uv$ both have poles
at $t=0$, yet $u^2-uv=1$. A counting path in
Lemma~\ref{lem:substitution} produces one monomial at a time. We need
a rule for combining its contributions that respects cancellation at
the chosen place.

We assign each weighted monomial $aX^\mu$ a matrix. For every regular
sum with a value in the prime field, the sum of its matrices acts by
that value on a single nonzero vector. In the example, the difference
of the two matrices acts as the identity on the vector. We use the size
parameter $M$ of Section~\ref{sec:descent} throughout.

\begin{lemma}[Matrices for the boundary value]\label{lem:pins}
Fix a singly exponential degree bound $L$. There are a nonzero vector
$\Omega\in\mathbb F_p^d$ and matrices $W_{a,\mu}\in\mathbb F_p^{d\times d}$
for $a\in E$ and $\mu\in\mathbb Z_{\ge0}^{h}$ with $|\mu|\le L$.
If $A=\sum_\mu a_\mu X^\mu\in E[X]$ has degree at most $L$, is
regular at the chosen place and has value $z\in\mathbb F_p$ there, then
\begin{equation}\label{eq:cartier-pin}
                 \biggl(\sum_\mu W_{a_\mu,\mu}\biggr)\Omega=z\Omega .
\end{equation}
In particular, this holds for $A=F(Y(X))$ of degree at most $L$,
with $F$ over $\mathbb F_p$ and $z=F(c)$. For fixed $\mu$ the matrix
is $\mathbb F_p$-linear in $a$, and $d$ is singly exponential. With
one polynomial-length advice string for all $a,\mu$ and all pairs of
indices, any specified entry is the number modulo $p$ of accepting
paths of a polynomial-time binary-guess computation.
\end{lemma}

The entries of $\Omega$ will be residues of test differentials on
the curve, followed by a field trace. If $A$ has boundary value
$z\in\mathbb F_p$, then $A^q-z=(A-z)^q$ vanishes to order at least
$q$ for every $p$-power $q$. For a test differential $\eta$ with pole
order less than $q$, it follows that $A^q\eta$ and $z\eta$ have the same
residue. The Cartier operator converts this observation into a rule
linear in $A$, through the identity $C_q(A^q\eta)=A C_q(\eta)$.
Speyer represented transitions of finite automata by related operators
on curve differentials. Bridy used this interpretation to bound the
number of states~\cite[Proposition~3.8 and Theorem~1.2]{Bridy}.
Here the tests must admit matrix entries that a counting path can
evaluate in the coordinates of the given curve equations.

\subsection{From the local coordinate to the curve equations}

Write $x=(x_1,\ldots,x_h)$ for the coordinate functions on the
geometric curve and $K=E(t,x)$ for the field they generate over $E$.
It is finite over $E(t)$. Indeed, each $x_i$ is algebraic over
$\overline{\mathbb F}_p(t)$, and only finitely many algebraic constants
occur in a defining equation for it. Restrict the chosen valuation to
$K$ and normalize its value group to $\mathbb Z$. Its residue field
$\kappa$ is finite, since the place lies over $t=0$ in a finite
function field extension. The residues of $Y(x)$ are still $c$.

Choose a uniformizer $s\in K$ at this place. The completion is
$\kappa((s))$, with the coefficient field obtained by lifting the
roots of $T^{|\kappa|}-T$. Put
\[
   J=\det\left(\frac{\partial g_i}{\partial X_j}\right),\qquad
   \eta_0=\frac{dt}{J(x,t)} .
\]
The Jacobian is nonzero at the generic point by
Proposition~\ref{prop:descent-curve}. Differentiating $g_i(x,t)=0$
expresses every $dx_j$ as a multiple of $dt$. In particular, $ds$
is a multiple of $dt$ because $s\in E(t,x)$. As $ds\ne0$ in
$\kappa((s))$, the differential $\eta_0$ is nonzero even when
$t$ is ramified at the branch.

Write $\operatorname{Res}\eta$ for the coefficient of $s^{-1}ds$ in
$\eta$. For $q=p^u$, $u\ge1$, Cartier acts in the local coordinate by
\[
 C_q\!\left(\sum_j v_j s^j\,ds\right)
       =\sum_i v_{q i+q-1}^{1/q}s^i\,ds .
\]
Grouping exponents modulo $q$ gives $C_q(D^q\eta)=D C_q(\eta)$,
and the coefficient at $i=-1$ gives
$\operatorname{Res}(C_q\eta)=(\operatorname{Res}\eta)^{1/q}$.
In the ambient coordinates $Z=(X_1,\ldots,X_h,t)$ put
\[
   \Lambda_q\!\left(\sum_j v_j Z^j\right)
        =\sum_i v_{q i+(q-1)\mathbf1}^{1/q} Z^i .
\]
This operator retains exponents congruent to $q-1$ in every
coordinate. The next formula computes the local action using only
this coefficient extraction and the defining equations of the curve.

\begin{lemma}\label{lem:cartier-formula}
For $f\in E[X_1,\ldots,X_h,t,t^{-1}]$ one has
\begin{equation}\label{eq:cartier-adjunction}
 C_q\bigl(f(x,t)\eta_0\bigr)
       =\Lambda_q\!\left(f\prod_{i=1}^h g_i^{q-1}\right)(x,t)\,\eta_0 .
\end{equation}
\end{lemma}

\begin{proof}
For a plane curve this is the formula of St\"ohr and
Voloch~\cite[Theorem~1.1]{StoehrVoloch}. The coordinate change in
higher codimension uses the naturality of Cartier on top
differentials~\cite{Cartier,Katz}. Recall its coefficient description.
In coordinates $z_1,\ldots,z_b$ forming a $p$-basis, each element
has a unique expansion in $z^j$, $0\le j_i<p$, over the subfield
of $p$-th powers. In a top form all terms with some $j_i<p-1$ have
an antiderivative in that coordinate. An exact form has zero
coefficient at $j=(p-1)\mathbf1$. Thus precisely this coefficient
survives modulo exact top forms. The intrinsic inverse Cartier map
from top forms to this quotient sends
\[
 a\,dv_1\wedge\cdots\wedge dv_b\longmapsto
 \left[a^p\bigwedge_{i=1}^b v_i^{p-1}dv_i\right].
\]
Indeed, $v^{p-1}dv$ is closed, obeys the product rule with scalars
raised to the $p$-th power, and is additive modulo exact forms.
The correction for additivity is the differential of the reduction
of the integral polynomial $((U+V)^p-U^p-V^p)/p$. This construction
commutes with the field embedding below, and inversion gives the
coefficient rule for Cartier. Iteration gives the rule for $q$.
In an iterated Laurent series field over a perfect field, grouping
exponents modulo $p$ shows that the series variables form a $p$-basis.

Introduce independent variables $y_1,\ldots,y_h$ over $\kappa((s))$.
Since $J(x,t)$ is invertible there, the formal implicit function
theorem gives $\widetilde x\in\kappa((s))[[y_1,\ldots,y_h]]^h$
with $\widetilde x(0)=x$ and $g_i(\widetilde x,t)=y_i$.
Hence $X\mapsto\widetilde x$, $t\mapsto t(s)$ embeds the ambient
rational function field in the iterated Laurent series field. Indeed,
the image contains the algebraically independent elements
$t(s),y_1,\ldots,y_h$. In the new variables the volume form is
\[
 dX_1\wedge\cdots\wedge dX_h\wedge dt
    = J(\widetilde x,t)^{-1}dy_1\wedge\cdots\wedge dy_h\wedge dt .
\]
Multiply by $f\prod_i g_i^{q-1}$ and apply Cartier. In the ambient
coordinates its coefficient becomes $\Lambda_q(f\prod_i g_i^{q-1})$.
In the series coordinates the factor $\prod_i y_i^{q-1}$ shifts
the constant term of $f(\widetilde x,t)/J(\widetilde x,t)$ to
$y^{(q-1)\mathbf1}$. This term alone contributes to the constant
$y$-coefficient after Cartier. Since $dt=t'(s)ds$, that coefficient
times $ds$ is $C_q(f(x,t)\eta_0)$. In the ambient calculation it is
$\Lambda_q(f\prod_i g_i^{q-1})(x,t)\eta_0$. This yields
\eqref{eq:cartier-adjunction}.
\end{proof}

\subsection{A finite family of test differentials}

We need a family large enough to detect a nonzero residue, whose span
is stable under each combined operation $\eta\mapsto a x^\mu C_q(\eta)$
with $|\mu|\le L$. Set $S=\sum_i\deg g_i$, $H=4(L+S+2)$, and take
the box of monomials
\[
  \mathcal B=\{X^\alpha t^j\mid 0\le\alpha_i\le H,\ -H\le j\le H\}.
\]
For $f\in\mathcal B$, $a,b\in E$ and $|\mu|\le L$, the Laurent polynomial
\begin{equation}\label{eq:cartier-row}
       aX^\mu\Lambda_q\!\left(bf\prod_i g_i^{q-1}\right)
\end{equation}
is supported in $\mathcal B$. In fact, the product with
$\prod_i g_i^{q-1}$ increases each upper exponent by at most
$(q-1)S$. After extraction and multiplication by $X^\mu$, the upper
bound in each coordinate is at most $H/q+S+L\le H$. The extracted
$X$ exponents are nonnegative, and the lower $t$ exponent is at least
$-H/q-1\ge-H$.

To bound pole orders of the tests, let $D$ be the degree of the
geometric curve. B\'ezout gives $D\le\prod_i\deg g_i$, which is
singly exponential by Proposition~\ref{prop:descent-curve}. On the
normalized projective model, a quotient of polynomials of degrees at
most $e$ has poles of order at most $eD$, provided its denominator
does not vanish identically on the curve. Homogenizing to a common
degree bounds the pole divisor by the intersection with the
denominator hypersurface. For $f/J$ one can take $e=(h+2)H+S$
after clearing negative powers of $t$. Restricting the valuation
to $K$ and normalizing only decreases this bound on pole orders.
Also $dt$ is regular since $t$ has positive valuation. Take $q$
to be the first positive power of $|E|$ exceeding $e\prod_i\deg g_i$.
This singly exponential $q$ works simultaneously for all the differentials
$f(x,t)\eta_0$ and fixes the scalars in $E$ under $q$-th powering.

Fix an $\mathbb F_p$-basis of $E$. Use pairs $(b,f)$ of a basis
element and a box monomial as our indices, and put
\[
 \ell(\eta)=\operatorname{Tr}_{\kappa/\mathbb F_p}(\operatorname{Res}\eta),
 \qquad \Omega(b,f)=\ell\bigl(bf(x,t)\eta_0\bigr).
\]
These entries are in $\mathbb F_p$, and $\ell(C_q\eta)=\ell(\eta)$
because the absolute trace is invariant under inverse Frobenius.

We show that the box contains a test with nonzero traced residue.
Choose a trace-one element of the finite separable extension
$\kappa/\mathbb F_p$ and lift it to an element $w$ of the valuation
ring of $K$. Then $\ell(w\,ds/s)=1$. The quotient
$(w\,ds/s)/\eta_0$ is in $K$. Express it as a finite linear combination
of monomials in $x$ with coefficients in $E(t)$, using the finiteness
of $E(t,x)/E(t)$. Approximate each rational coefficient by a Laurent
polynomial in $t$ over $E$. Since $t$ has positive valuation, sufficiently
long approximations preserve the residue after multiplication by the
finitely many test factors $x^\alpha\eta_0$. This gives a Laurent
polynomial $f_0\in E[X,t,t^{-1}]$ with $\ell(f_0(x,t)\eta_0)=1$.

The exponents of $f_0$ need not yet lie in $\mathcal B$. Apply
Lemma~\ref{lem:cartier-formula} with a $p$-power $Q$ larger than
the absolute values of all these exponents. The polynomial
$\Lambda_Q(f_0\prod_i g_i^{Q-1})$ is supported in $\mathcal B$ by the
same exponent bound, and its traced residue against $\eta_0$ is still
one. Expand its coefficients in the chosen basis of $E$. At least
one of the resulting entries $\Omega(b,f)$ is nonzero. The power $Q$
and the local functions $w,s$ serve only to prove this existence.
The matrix entries below use the uniform power $q$ and the curve equations.

\subsection{The matrix action and its entry counts}

\begin{proof}[Proof of Lemma~\ref{lem:pins}]
In row $(b,f)$ of $W_{a,\mu}$, enter the coefficients of
\eqref{eq:cartier-row} in the basis $b'f'$, with columns indexed
by $(b',f')$. The support bound makes this a square matrix over
$\mathbb F_p$. Multiplication by $a$ after extraction makes the
row $\mathbb F_p$-linear in $a$. By
Lemma~\ref{lem:cartier-formula}, the $(b,f)$ entry on the left of
\eqref{eq:cartier-pin}, with evaluations on the curve understood, is
\[
   \ell\!\left(A\,C_q(bf\eta_0)\right)
    =\ell\!\left(C_q(A^q bf\eta_0)\right)
    =\ell(A^q bf\eta_0)=z\,\ell(bf\eta_0).
\]
The last equality uses regularity of the whole sum $A$. The zero
of $A^q-z$ has order at least $q$, so $(A^q-z)bf\eta_0$ has zero
residue. This proves \eqref{eq:cartier-pin}. More generally, if
the boundary value is $z\in E$, then $z^q=z$ and the right side
becomes $R_z\Omega$, where row $(b,f)$ of $R_z$ expands $zb$ in
the fixed basis while keeping $f$ fixed.

There are $d=[E:\mathbb F_p](H+1)^h(2H+1)$ indices, each of
polynomial bit length. To compute an entry in row $(b,f)$ and
column $(b',f')$, subtract $(\mu,0)$ from the exponent vector of
$f'$ and call the difference $i$. With
\[
       j=qi+(q-1)\mathbf1,\qquad V=bf\prod_i g_i^{q-1},
\]
the requested entry is the $b'$-coordinate of $a[Z^j]V$. The
inverse Frobenius in $\Lambda_q$ fixes $E$ since $q$ is a power
of $|E|$. If $j$ lies outside the exponent vector of $f$ plus
the box $[0,(q-1)S]^{h+1}$, the entry is zero.

Otherwise take the finite extension $E'/E$ of least positive degree
such that $N=|E'|-1>(q-1)S$. Its elements have polynomial bit length,
since $|E'|\le |E|((q-1)S+2)$. Orthogonality in $(E')^\times$
gives
\[
       [Z^j]V=N^{-(h+1)}
           \sum_{z\in((E')^\times)^{h+1}} V(z)z^{-j}.
\]
Every coordinate interval of the support has width less than $N$.
Thus the only exponent of the support that can be congruent to $j$
modulo $N$ in all coordinates is $j$ itself. The argument applies
also to the negative powers of $t$.

Give the counting procedure the curve evaluators from
Proposition~\ref{prop:descent-curve}, the exponent defining $q$, field
presentations, embeddings, and basis data as polynomial-length advice. A path guesses one tuple
of nonzero elements of $E'$ in canonical encodings, rejecting invalid
guesses. It evaluates $V(z)z^{-j}$ and the normalization factor by
the curve evaluators and binary powering. Multiply by $a$ and apply
the coordinate projection for $b'$, extended $\mathbb F_p$-linearly
from $E$ to $E'$. Both operations take polynomial time and commute
with the sum over tuples. If the resulting prime-field element has
integer representative $v\in[0,p-1]$, a final guess of
$\lceil\log_2 p\rceil$ bits accepts exactly when its value is less
than $v$. Each summand therefore contributes the required number
of accepting paths modulo $p$.

We may pad $d$ to the next power of two by adding zero rows and
columns to the matrices and zero coordinates to $\Omega$. The
entry procedure then takes binary indices of one fixed length.
Section~\ref{sec:border} uses one such index at which $\Omega$ is
nonzero and the entry procedure common to all the monomials.
\end{proof}
